\subsection{6.2 Lanczos zählt Hamiltonschritte, nicht
Bosonen}\label{lanczos-zuxe4hlt-hamiltonschritte-nicht-bosonen}

Die einfache Kette entsteht aus \(F,HF,H^2F,\ldots\) durch
Orthogonalisierung:

\[Hu_n=b_nu_{n-1}+a_nu_n+b_{n+1}u_{n+1}.\]

Die aktuellen Koeffizienten sind

\[\frac{(a_0,a_1,a_2,a_3,a_4)}\Delta
=\left(0,1,2,3,\frac{105168998}{26292551}\right),\]

\[\frac{(b_1^2,b_2^2,b_3^2,b_4^2)}{g^2}
=\left(480,916,\frac{299520}{229},\frac{78877653}{47632}\right).\]

Der fünfte Vektor enthält sowohl \(v_4\) als auch \(w_2\). Er hat keine
feste Bosonzahl. Eine Kettenlänge von höchstens 33 folgt deshalb nicht
aus den 33 möglichen Bosonzahlen. Der Lanczosindex ist ein Rechenindex;
seine Deutung als räumlicher Ort wäre ein zusätzlicher unbegründeter
Schritt.

Die fünf Glieder bestimmen die zehn Momente
\(\mu_n=\langle F,H^nF\rangle\) für \(0\le n\le9\). Ein Weg zum
ausgelassenen sechsten Kettenvektor und zurück braucht mindestens zehn
Schritte. Insbesondere

\[\mu_2=480g^2,\quad \mu_3=480\Delta g^2,\quad
\mu_4=480\Delta^2g^2+670080g^4.\]

Diese Momente gehören \(F\), nicht \(\Omega\) und nicht den geladenen
Antwortoperatoren. Die rationale Sturm-Einschließung der kleinsten
Ritzenergie

\[-1.12963813<E_{\mathrm{Ritz},5}/\Delta<-1.12963811\]

wurde hier unabhängig aus den gegebenen Koeffizienten nachvollzogen. Der
Ursprung des übernommenen vierten Bosonmoments wurde nicht neu
enumeriert. {[}S068, §4{]}

\subsection{6.3 Ein exakter Nachweis fehlender
Auswahl}\label{ein-exakter-nachweis-fehlender-auswahl}

Aus den Momenten folgt

\[\frac{\mu_4\mu_2}{\mu_3^2}=1+1396(g/\Delta)^2.\]

Das ergibt \(4.49\) bei \(g/\Delta=1/20\), aber \(1.8725\) bei \(1/40\).
Beide Modelle besitzen dasselbe \(W\), dieselben inneren Symmetrien und
denselben kommutierenden Clock. Die dimensionslose Differenz ist keine
Einheitenfrage. Damit ist belegt, dass diese Grunddaten allein die
physikalische Dynamik nicht eindeutig auswählen. Die Referenz \(F\) muss
für eine tatsächliche Messung verfügbar sein; die algebraische
Ungleichheit der Prozesse ist davon getrennt. {[}S068, §8{]}

\subsection{6.4 Nachgereichter Text: die interessante Fortsetzung liegt
in
Matrixmomenten}\label{nachgereichter-text-die-interessante-fortsetzung-liegt-in-matrixmomenten}

Der während dieser Arbeit nachgereichte Text \emph{Operationssatz,
nativer Grundzustand und Feldwörterbuch}, beschriftet als v1.6.3, ist
byteidentisch mit dem bereits gelesenen aktuellen Repositorybericht
derselben Linie. Sein Tensor, Lochbild, Kommutantenstufen und
Nichtschlussbefund sind deshalb bereits berücksichtigt. Seine Priorität
„den Zwei-Boson-Singulettraum ausschreiben`` ist in v1.6.7 inzwischen
durch die vier expliziten Richtungen erfüllt. Die pauschale Behauptung,
der nichtskalare Clock liege außerhalb der zusammenhängenden Gruppe, ist
durch das konkrete Spin-Wort überholt. {[}S073; S033, §3; S068, §3{]}

\textbf{Eine weiterhin interessante Richtung steckt in seinen
Onishi-/Wick-Spurnetzwerken.} Sie berechnen hohe Bosonstufennormen, ohne
sämtliche Fockamplituden auszuschreiben. Der Nichtschluss beweist
jedoch, dass die skalaren Normen \(\nu_k\) allein nicht genügen. Die
natürliche Weiterentwicklung sind gemischte Wortmomente und
projektoraufgelöste Matrixmomente, beispielsweise

\[\Gamma_{(R,u),(R',v)}
=\langle uP_Rv_2,vP_{R'}v_2\rangle.\]

Hier sind \(u,v\) ausdrücklich Wörter in \(H\), \(T_+\), \(T_-\) und den
benötigten tatsächlichen Observablen. Solche Matrizen behalten
Kreuzterme zwischen den vier Richtungen und nachfolgenden Bosonstufen.
Sie passen genau zum Prozesskern aus Abschnitt 8. Ein tensorbasierter
Kontraktionsalgorithmus dafür könnte die gemeinsame Rechnung wesentlich
verkleinern, ohne den verzweigten Raum irrtümlich zu einer einzigen
Bosonleiter zu machen. Die Entwicklung und Ausführung dieses erweiterten
Kontraktionsalgorithmus ist in diesem Paper nicht bereits erfolgt.

Eine dazu passende \textbf{exakte analytische Fortsetzung} ist die
Blockresolvente. Für eine orthogonale Zerlegung und

\[H=\begin{pmatrix}A&B^\dagger\\B&D\end{pmatrix}\]

gilt für \(\operatorname{Im}z\ne0\)

\[P(z-H)^{-1}P
=\big[z-A-B^\dagger(z-D)^{-1}B\big]^{-1}.\]

Dies folgt durch Eliminieren der zweiten Blockzeile. Der Rest
verschwindet nicht; er erscheint als energieabhängige Selbstenergie. Für
\(\eta=|\operatorname{Im}z|>0\) und endliche selbstadjungierte Blöcke
folgt

\[\|\Sigma(z)\|\le\|B\|^2/\eta,\qquad
\|P(z-H)^{-1}P-(z-A)^{-1}\|\le\|B\|^2/\eta^3.\]

Die erste Schranke benutzt \(\|(z-D)^{-1}\|\le1/\eta\), die zweite die
Resolventenidentität und die Schranken \(1/\eta\) für beide äußeren
Resolventen. Diese Kontrolle wird nahe der reellen Achse schwach;
präzise Polaussagen benötigen eine bessere Rest- und Abstandskontrolle.
Sie ist trotzdem ein klarer Erfolgstest für die vorgeschlagene
Spurnetzwerkfortsetzung: nicht nur weitere Normen berechnen, sondern die
tatsächliche Matrixrestkopplung der gemeinsamen Antwort einschließen.

Die Feldgradierungen des Textes bleiben ebenfalls als Kandidaten
relevant. Die dort genannten Aufteilungen 24/36 oder 40/20 sind
ansatzgebundene Vektor-/Tensorkanäle, keine Auswahl eines gesunden
relativistischen Parents. Der nächste Gewinn wäre ein vollständiger
positiver Feldvertrag für genau einen solchen Ansatz mit nachgewiesener
Rückführung auf W. Die Aussage „einziger einheitlicher Feldtyp`` gilt
nur in der abgegrenzten Ein-Kopie-, Händigkeit- und Ableitungsklasse.

\section{7. Präparation, Restantwort und die Herkunft von
Teilen}\label{pruxe4paration-restantwort-und-die-herkunft-von-teilen}

\subsection{7.1 Ein kontrollierter Weg in den nativen
Grundzustand}\label{ein-kontrollierter-weg-in-den-nativen-grundzustand}

Die zweite v1.6.6-Linie betrachtet

\[H(s)=\Delta N_b+sgX,\quad0\le s\le1\]

im festen \(N=64\)-Sektor. Der Startzustand \(F\) ist dort bei \(s=0\)
eindeutig. Aus einer Komplementuntergrenze und Variationsobergrenzen
folgt eine durchgehende Lücke

\[\operatorname{gap}H(s)>
\frac{97697}{1187646}\Delta>0.0822610441\Delta.\]

Mit einer konservativen quantitativen Adiabatenschranke erhält man bei
\(T\Delta=10^9\) einen Fehlerbetrag unter \(0.000523\) und eine
Grundzustandswahrscheinlichkeit über \(99.9999726\%\). Das ist ein
mathematisch bedingter Präparationssatz unter steuerbarer Rampe und
verfügbarem \(F\), keine durchgeführte Rampensimulation oder
Laborpräparation. Die globale Nullenergieentartung bei \(s=0\) stört
diesen festen Sektor nicht, solange die Ladung exakt erhalten bleibt.
{[}S001, §2; Adiabatensatz E3{]}

\subsection{7.2 Die volle Antwort ist nicht exakt
zweilinienförmig}\label{die-volle-antwort-ist-nicht-exakt-zweilinienfuxf6rmig}

Mit \(b=\langle N_b\rangle\) und

\[S=15-\frac7{32}b,\quad a=\frac{\Delta b-E_0}{64}\]

lauten die bisherigen normierten Spektralmomente

\[m_0=1,\quad m_1=0,\quad m_2=g^2S,\quad
m_3=g^2(\Delta S+7a).\]

Für den zweiten orthogonalen Operatorrest folgt

\[\rho^2=m_4-\frac{m_3^2}{m_2}-m_2^2,\qquad
\beta_2^2=\rho^2/m_2.\]

Die neue Restidentität liefert unter den nativen Voraussetzungen

\[\operatorname{Var}(N_b)>0.1193989,
\qquad\beta_2^2>9.640875\cdot10^{-7}\Delta^2.\]

Der Koeffizient ist eine Kopplung an den restlichen Antwortbereich. Er
ist weder dessen prozentuales Gewicht noch eine punktweise Schranke der
vollständigen Selbstenergie. Eine brauchbare Näherung benötigt außerdem
eine obere Fehlerkontrolle im gewünschten Frequenzfenster. {[}S001,
§3{]}

\subsection{7.3 Eine Bank ist in ihrer ursprünglichen Modenbasis stark
verbunden}\label{eine-bank-ist-in-ihrer-urspruxfcnglichen-modenbasis-stark-verbunden}

Die frische Rechnung am gepinnten \(W\) ergibt 480 Kanten auf 64
Fermionmoden, Grad 15 und das exakte Adjazenzspektrum

\[15^{[1]},\quad3^{[25]},\quad(-1)^{[30]},\quad(-5)^{[3]},\quad(-9)^{[5]}.\]

Die Laplacelücke ist zwölf. Für eine Teilmenge mit \(s\) Moden gilt
deshalb

\[|\partial S|\ge \frac{12s(64-s)}{64}
=\frac{3s(64-s)}{16}.\]

Ein 32/32-Schnitt kreuzt mindestens 192 Paarterme. Der Beweis nutzt den
zentrierten Indikator \(1_S-(s/64)\mathbf1\): Sein Laplacequadratikwert
zählt die Schnittkanten und seine Norm ist \(s(64-s)/64\).

Das erklärt, weshalb das bloße Aufteilen der ursprünglichen Koordinaten
keine zwei unabhängigen Banken ergibt. Es beweist keine Raumdimension
und schließt anders konstruierte emergente Lokalität nicht aus. Eine
neue lokale Basis müsste Operatoren, Zustand und Auslesung gemeinsam
transformieren. {[}S001, §4; eigene Tensorprüfung{]}

\section{8. Eigene Verbindung I: ein gemeinsamer
Prozesskern}\label{eigene-verbindung-i-ein-gemeinsamer-prozesskern}

\subsection{8.1 Die zentrale
Vereinfachung}\label{die-zentrale-vereinfachung}

Die bisherigen Teilprobleme können auf ein gemeinsames mathematisches
Objekt bezogen werden: die Überlappungen aller tatsächlich zulässigen
Operationsgeschichten. Sei \(\mathcal W\) die Menge endlicher Wörter aus
einem ausdrücklich festgelegten Alphabet von Operationen und ihren
Adjungierten. Auf einem gemeinsamen Startzustand \(\Omega\) definiere

\[\boxed{\Gamma(u,v)=\langle u\Omega,v\Omega\rangle
=\omega(u^\dagger v).}\]

Dieser \textbf{Prozesskern} umfasst nicht allein Wahrscheinlichkeiten
einzelner Ausgänge. Seine komplexen Kreuzterme halten fest, welche
Geschichten noch interferieren können. Werden Register und Apparate in
den gemeinsamen Raum aufgenommen, gehören auch ihre Zustandsänderungen
zu den Wörtern.

Die Idee ist mathematisch mit der GNS-Rekonstruktion positiver
Funktionale verwandt. Für gemischte Anfangszustände kann man entweder
das Funktional direkt verwenden oder eine gemeinsame Reinigung mitsamt
unzugänglicher Umgebung führen. Die hier ausgeschriebene Anwendung ist
ein Verbindungssatz für die TFPT-Arbeitsstände, kein literaturweiter
Neuheitsanspruch.

\subsection{8.2 Satz: Rekonstruktion aus einem vollständigen positiven
Wortkern}\label{satz-rekonstruktion-aus-einem-vollstuxe4ndigen-positiven-wortkern}

\textbf{Voraussetzungen.} Gegeben seien eine unital abgeschlossene
\(*\)-Algebra \(\mathfrak A\) der zulässigen Wörter, ihre Relationen und
ein normiertes positives Funktional \(\omega\) mit \(\omega(1)=1\). Für
eine Darstellung durch beschränkte Operationen werde zusätzlich für
jeden Generator \(a\) eine Schranke

\[\omega(x^\dagger a^\dagger ax)\le C_a^2\omega(x^\dagger x)
\qquad(x\in\mathfrak A)\]

verlangt. In endlichdimensionalen konkreten Modellen ist diese Bedingung
erfüllt. Ein unbeschränkter Hamiltonoperator benötigt außerhalb dieser
Fassung eigene Definitionsbereichs- und Selbstadjungiertheitskontrolle.

\textbf{Aussage.} Es gibt einen Hilbertraum \(\mathcal H_\omega\), einen
zyklischen Vektor \(\Omega_\omega\) und eine Darstellung \(\pi\) mit

\[\omega(a)=\langle\Omega_\omega,\pi(a)\Omega_\omega\rangle.\]

Auf dem von allen Wörtern erzeugten Raum ist diese Darstellung bis auf
unitäre Äquivalenz eindeutig.

\textbf{Beweis.} Auf \(\mathfrak A\) setze man
\(\langle x,y\rangle=\omega(x^\dagger y)\) und quotiente die
Nullvektoren aus. Positivität liefert Cauchy--Schwarz. Die
Generatorenschranke zeigt, dass ein Nullvektor nach Linksanwendung von
\(a\) wieder ein Nullvektor ist. Also ist \(\pi(a)[x]=[ax]\)
wohldefiniert und beschränkt. Die Vervollständigung ergibt den
Hilbertraum; \([1]\) ist zyklisch. Zwei solche Realisierungen haben auf
allen endlichen Linearkombinationen derselben Wörter dieselben
Skalarprodukte. Die Zuordnung \(u\Omega\mapsto u\widetilde\Omega\) ist
daher wohldefiniert und isometrisch und erweitert sich wegen Zyklizität
zu einer unitären Abbildung. \(\square\)

Der Satz erklärt, in welchem Sinn ein „Universalraum`` als gemeinsamer
erreichbarer Prozessraum präzise sein kann. Er setzt jedoch den
vollständigen positiven Kern und das Operationsalphabet voraus. Eine
einzelne endliche Momentenmatrix genügt weder für eine eindeutige
unendliche Fortsetzung noch für physikalische Lokalität.

Für ein physikalisches Instrument müssen zusätzlich die
Ausgangsoperatoren \(M_j\) eine vollständige Operation bilden, etwa
\(\sum_jM_j^\dagger M_j=I\) für einen spurtreuen Gesamtprozess. Die
Wahrscheinlichkeiten sind dann \(\omega(M_j^\dagger M_j)\); selektive
Zweige behalten ihre tatsächliche Ausbeute. Positivität allein setzt
diese Normierungs- und Kompositionsrelationen nicht ein.

\textbf{Mehrere Schatten derselben Quelle.} Bei festgelegten linearen
Reduktionen \(R_i\) ist eine endliche uneingeschränkte Zustandsklasse
genau dann gemeinsam unterscheidbar, wenn auf hermiteschen spurlosen
Differenzen \(\bigcap_i\ker R_i=\{0\}\) gilt. Alle echten Teilmarginalen
von \((|000\rangle\pm|111\rangle)/\sqrt2\) sind dagegen gleich; erst
eine phasensensitive gemeinsame Messung wie \(XXX\) unterscheidet sie.
Auch passende paarweise Überschneidungen garantieren noch keine
gemeinsame globale Realisierung. Das Wortkernkriterium verlangt daher
gemeinsame zulässige Operationen einschließlich der relevanten
Kreuzterme. Es rekonstruiert den zyklisch erreichbaren Raum, keinen
völlig entkoppelten dunklen Zusatzsektor. {[}S074, B5.3{]}

\subsection{8.3 Zustandswahl und Dynamik lassen sich nicht durch
Rekonstruktion
erschleichen}\label{zustandswahl-und-dynamik-lassen-sich-nicht-durch-rekonstruktion-erschleichen}

Der Satz konstruiert aus vorhandenen Prozessdaten eine gemeinsame
Darstellung. Er erzeugt diese Daten nicht aus E8 allein. Die Familie
\(H_{\Delta,g}\) aus Abschnitt 6 zeigt ausdrücklich, dass gleiche
Grammatik verschiedene positive Prozesskerne zulässt. Ebenso können
verschiedene Anfangszustände auf derselben Algebra verschiedene
Funktionale liefern.

Für eine kontinuierliche Zeitentwicklung benötigt man zusätzlich eine
konsistente, ausreichend stetige Automorphismengruppe \(\alpha_t\). In
einer invariant gewählten GNS-Darstellung kann sie unter den üblichen
Voraussetzungen unitär implementiert werden. Eine einparametrige Gruppe
und ihre Stetigkeit fallen nicht aus einer Ordnung-6-Clock heraus. Der
Energie-Nullpunkt und die positive Energiebedingung müssen ebenfalls
korrekt behandelt werden.

Der wissenschaftliche Gewinn liegt deshalb in einer scharfen Reduktion:
Die gemeinsame Herkunft wird zur Auswahl eines \textbf{positiven,
zeitlich konsistenten, lokal komponierbaren Prozesskerns}. Das ist ein
präziseres Ziel als eine Liste voneinander unabhängiger schöner
Matrizen.

\subsection{8.4 Kompatibilitätssatz für verschiedene
Arbeitsstände}\label{kompatibilituxe4tssatz-fuxfcr-verschiedene-arbeitsstuxe4nde}

Angenommen, zwei Konstruktionen verwenden dasselbe Wortalphabet und
besitzen für \textbf{alle} zulässigen Wörter denselben Gramkern. Dann
sind ihre zyklischen Ausführungen unitär äquivalent und die
entsprechende Unitäre verschränkt sämtliche Generatoren. Gleichheit nur
von \(\Gamma(1,u)\) oder wenigen Wahrscheinlichkeiten genügt nicht.

Der lokale Recordsatz ist der einfachste Spezialfall. Ist
\(K^\dagger K=L^\dagger L\), definiert man \(V(Kx)=Lx\). Aus \(Kx=0\)
folgt \(\|Lx\|^2=0\), also ist die Definition wohldefiniert. Sie erhält
Skalarprodukte und ist auf dem Bild isometrisch. Umgekehrt ergibt
\(L=VK\) unmittelbar die Gramgleichheit.

Für ganze Prozesse lautet die entsprechende Prüffrage:

\[\langle u\Omega,v\Omega\rangle
\stackrel{?}=\langle\widetilde u\widetilde\Omega,
\widetilde v\widetilde\Omega\rangle\]

für dieselben Eingriffe, Referenzen und Ausgänge. Damit wird verhindert,
dass ein N=9-Transferzeuge, eine N=4-Zellreferenz und die
N=64-Grundzustandsantwort ohne Adapter als eine einzige Lösung gelten.

\subsection{8.5 Ein endlicher Fehlervertrag für
Adapter}\label{ein-endlicher-fehlervertrag-fuxfcr-adapter}

Sei \(J:\mathcal H_{\mathrm{eff}}\to\mathcal H\) eine Isometrie,
\(h=J^\dagger HJ\) und

\[R=HJ-Jh=(I-JJ^\dagger)HJ.\]

Für endliche selbstadjungierte Matrizen gilt durch die Duhamel-Identität

\[\boxed{\|e^{-itH}J-Je^{-ith}\|\le |t|\,\|R\|.}\]

\textbf{Beweis.} Man differenziert \(e^{-i(t-s)H}Je^{-ish}\) nach \(s\),
integriert von null bis \(t\) und benutzt die Norm eins beider unitären
Faktoren. Für negatives \(t\) folgt dieselbe Schranke mit \(|t|\).
\(\square\)

Bei einem Effekt \(0\le E\le I\) ist der Unterschied der
Ausgangswahrscheinlichkeiten zweier so verglichener normierter Zustände
höchstens \(2|t|\|R\|\). Damit erhält das Wort „gemeinsam`` einen
zahlenmäßigen Inhalt: Eine behauptete effektive Wirkung muss größer als
die kontrollierte Adapterunsicherheit sein. Bei langen Zeiten kann die
lineare Schranke unbrauchbar werden; dann benötigt man bessere spektrale
oder lokale Abschätzungen. Die Methode ersetzt insbesondere keine
uniforme Kontrolle im thermodynamischen Limes.

\section{9. Eigene Verbindung II: Transfer mit mitgerechneter
Ladungsreferenz}\label{eigene-verbindung-ii-transfer-mit-mitgerechneter-ladungsreferenz}

